%# -*- coding:utf-8 -*-
%% Based on moderncv's template_en.tex.
%% Copyright 2006-2008 Xavier Danaux (xdanaux@gmail.com).
%
% This work may be distributed and/or modified under the
% conditions of the LaTeX Project Public License version 1.3c,
% available at http://www.latex-project.org/lppl/.

\documentclass[11pt,a4paper]{moderncv}

\usepackage{fontspec,xunicode}
\usepackage[slantfont,boldfont]{xeCJK}
\usepackage{xcolor}

% PATCH (latextemplates.com): upstream picked fonts with \ifwindows/\ifmacosx
% branches (Times New Roman, SimSun/STSong/STFangsong — all OS-bundled fonts
% looked up by fontconfig FAMILY name). None of those branches match Linux,
% and this pinned image's fontconfig only resolves DejaVu and the URW
% base-35 clones anyway (see research/README.md's font-trap section) — the
% upstream fonts would silently fail to load with a clean-looking compile.
% Replaced with TeX Gyre Termes (a Times-alike) for the Latin text and the
% Fandol family for CJK, both loaded by FILENAME via kpathsea, which this
% TeX Live resolves regardless of platform or fontconfig.
\setmainfont{texgyretermes-regular.otf}[
  BoldFont = texgyretermes-bold.otf,
  ItalicFont = texgyretermes-italic.otf,
]
\setCJKmainfont[ItalicFont=FandolKai-Regular.otf, BoldFont=FandolHei-Bold.otf]{FandolSong-Regular.otf}
\setCJKsansfont{FandolHei-Regular.otf}
\setCJKmonofont{FandolSong-Regular.otf}

% moderncv themes
\moderncvtheme[blue]{classic}  % optional argument are 'blue' (default), 'orange', 'red', 'green', 'grey' and 'roman' (for roman fonts, instead of sans serif fonts)

% adjust the page margins
\usepackage[scale=0.9]{geometry}
\AtBeginDocument{\recomputelengths}  % required when changes are made to page layout lengths

% personal data
\firstname{Wei}
\familyname{Chen}
\title{Software Engineer / 软件工程师}
\address{Hangzhou, China}{}
\mobile{+86 138 0000 0000}
\email{wei.chen@example.com}
\homepage{wei-chen.dev}
\social[github]{weichen-dev}
\extrainfo{%
  LinkedIn: linkedin.com/in/weichen-dev
}

% \photo[64pt]{avatar.png}  % add a photo by dropping a file next to main.tex and uncommenting this line

%----------------------------------------------------------------------------------
%            content
%----------------------------------------------------------------------------------
\begin{document}
\maketitle
\vspace*{-14mm}

\section{工作经历 / Experience}
\cventry{2022.03--至今}{高级后端工程师}{字节跳动}{北京}{}{负责推荐系统服务端架构，
  将核心排序服务的 P99 延迟从 180ms 降到 60ms，主导了从单体服务到微服务的拆分}
\cventry{2019.07--2022.02}{后端工程师}{美团}{北京}{}{参与外卖订单系统的开发与维护，
  设计并实现了订单状态机，支撑日均千万级订单}
\cventry{2018.01--2019.06}{实习生}{腾讯}{深圳}{}{参与微信支付对账系统开发，
  独立完成了对账差异自动修复模块}

\section{教育经历 / Education}
\cventry{2015.09--2019.06}{学士}{浙江大学}{计算机科学与技术}{}{}
\cvlistitem{校级一等奖学金 (2017, 2018)}
\cvlistitem{ACM/ICPC 区域赛银奖}

\section{项目 / Projects}
\subsection{开源项目 / Open Source}
\cvline{fastcache}{一个用 Go 编写的高性能本地缓存库
  \emph{github.com/weichen-dev/fastcache} \textbf{2,300+ stars}}
\cvline{ordergraph}{订单状态可视化调试工具，被团队内部广泛使用}

\section{技能 / Skills}
\cvline{\textbf{后端}}{Go, Java, PostgreSQL, Redis, Kafka}
\cvline{\textbf{基础设施}}{Kubernetes, Docker, Terraform}
\cvline{\textbf{其他}}{熟练使用 Git/Vim, 英语工作沟通流利}

\section{Publications}
\cvline{2021}{W. Chen, et al. \emph{Reducing Tail Latency in Large-Scale
  Recommendation Serving.} Internal Engineering Report.}

\end{document}

%% end of file `main.tex'.

%%% Local Variables:
%%% mode: latex
%%% TeX-command-extra-options: "-shell-escape"
%%% TeX-master: t
%%% TeX-engine: xetex
%%% End:
